\chapter{Fixed-Damping Spectral Theory}

The first level of the hierarchy assumes that damping is a fixed matrix. This
is the last level at which one can still try to extract a large amount of
engineering information from the graph Laplacian \(L\), the inertia matrix
\(M\), and the damping matrix \(D\) without retaining converter-control
states. The chapter develops that analysis aggressively. It derives the
spectral screens, the exact first-order eigenvalue sensitivities, the
decoupled modal formulas for \(d/dL\), \(d/dD\), and \(d/dM\), and the limits
of graph-only reinforcement.

The conclusion is deliberately sharp. With fixed \(D\), adding inertia is not
automatically stabilizing for oscillatory damping margins. Adding a line is not
automatically stabilizing either. Damping ratio is a quotient: it improves when
modal damping rises faster than modal frequency. If an action raises stiffness
or inertia but does not raise effective modal damping, the damping ratio can
fall even while a graph metric improves.

\section{The Classical Swing Operator}

Consider
\begin{equation}
  M\ddot\theta+D\dot\theta+L\theta=p,
  \label{eq:fixed_swing}
\end{equation}
with \(M\succ0\), \(D\succeq0\), and \(L=L^T\succeq0\) on the projected
subspace orthogonal to uniform angle shift. The homogeneous modal ansatz
\(\theta(t)=ve^{st}\) gives the QEP
\begin{equation}
  P(s)v=(s^2M+sD+L)v=0.
  \label{eq:fixed_qep}
\end{equation}
This is the basic fixed-\(D\) operator. All planning screens in this chapter
are approximations or consequences of \eqref{eq:fixed_qep}.

In inertia-normalized coordinates,
\begin{equation}
  \ddot\vartheta+\Dt\dot\vartheta+\Lt\vartheta=0,
  \label{eq:fixed_normalized}
\end{equation}
where
\begin{equation}
  \Lt=M^{-1/2}LM^{-1/2},\qquad
  \Dt=M^{-1/2}DM^{-1/2}.
\end{equation}
Let \(\Lt Q=Q\Lambda\). In modal coordinates \(z=Q^T\vartheta\),
\begin{equation}
  \ddot z+\Gamma \dot z+\Lambda z=0,\qquad
  \Gamma=Q^T\Dt Q.
  \label{eq:fixed_modal}
\end{equation}
The graph diagonalizes stiffness. It diagonalizes damping only when
\begin{equation}
  \Lt\Dt=\Dt\Lt
  \label{eq:commutation_condition}
\end{equation}
on the nonzero modal subspace. Equation~\eqref{eq:commutation_condition} is the
first hidden assumption behind graph-only damping arguments.

A basis-dependent off-diagonal ratio is useful only when the eigenvalues of
\(\Lt\) are well separated. For clustered eigenvalues, the basis inside the
cluster can rotate without changing the physics. A coordinate-free global gate is
\begin{equation}
  \eta_{\rm comm}
  =
  \frac{\norm{\Lt\Dt-\Dt\Lt}_F}
  {2\norm{\Lt}_F\norm{\Dt}_F+\varepsilon_{\rm reg}}.
  \label{eq:commutator_nonprop_index}
\end{equation}
For two spectral clusters \(I\) and \(J\), with orthogonal projectors
\(P_I\) and \(P_J\), the cross-cluster damping gate is
\begin{equation}
  \eta_{IJ}
  =
  \frac{\norm{P_I\Dt P_J}_2}
  {\operatorname{sep}(\Lambda_I,\Lambda_J)+\varepsilon_{\rm reg}}.
  \label{eq:projector_cluster_nonprop}
\end{equation}
This form is invariant to rotations of the modal basis inside each cluster and
is the safer trigger for adaptive reduced-QEP construction.

\section{Exact QEP Energy Identities}

The scalar decoupled formulas used later as screens have an exact ancestor in
the full QEP. Let \(P(s)x=0\) for
\[
  P(s)=s^2M+sD+L,
  \qquad
  M\succ0,\quad D\succeq0,\quad L\succ0
\]
on the nonzero-angle subspace. Assume \(s=\sigma+j\omega\) with
\(\omega\ne0\). Multiplying by \(x^H\) gives
\begin{equation}
  s^2 m+s d+\ell=0,
  \qquad
  m=x^HMx,\quad d=x^HDx,\quad \ell=x^HLx.
  \label{eq:qep_scalar_energy_identity}
\end{equation}
Since \(M,D,L\) are Hermitian in this fixed-\(D\) layer, \(m,d,\ell\) are real
and \(m>0\). Separating real and imaginary parts yields
\begin{align}
  \sigma
  &=
  -\frac{x^HDx}{2x^HMx}, \label{eq:qep_exact_sigma}\\
  |s|^2
  &=
  \frac{x^HLx}{x^HMx}, \label{eq:qep_exact_modulus}\\
  \zeta
  &=
  \frac{x^HDx}
  {2\sqrt{(x^HMx)(x^HLx)}}. \label{eq:qep_exact_zeta}
\end{align}
These identities do not require \(D=\gamma M\). They require the fixed
Hermitian QEP structure and an oscillatory root.

\begin{proposition}[Absence of undamped oscillatory modes]
For the fixed Hermitian QEP, an oscillatory eigenvalue on the imaginary axis
\(s=j\omega\), \(\omega\ne0\), exists only if there is a nonzero vector \(x\)
such that
\begin{equation}
  Dx=0,\qquad Lx=\omega^2Mx.
  \label{eq:undamped_mode_condition}
\end{equation}
\end{proposition}

\begin{proof}
If \(s=j\omega\), then \eqref{eq:qep_exact_sigma} gives \(x^HDx=0\). Since
\(D\succeq0\), this implies \(Dx=0\). The QEP then reduces to
\((-\,\omega^2M+L)x=0\). The converse follows by substitution.
\end{proof}

\section{Generalized Graph Modes}

For placement calculations it is often cleaner to use the generalized
eigenproblem
\begin{equation}
  L\phi_k=\nu_k M\phi_k,\qquad
  \phi_k^T M\phi_\ell=\delta_{k\ell}.
  \label{eq:generalized_modes}
\end{equation}
The zero mode is removed. The scalar modal stiffness is \(\nu_k\), and the
frozen-mode scalar damping is
\begin{equation}
  \delta_k=\phi_k^TD\phi_k.
  \label{eq:modal_delta_physical}
\end{equation}
If the modes are decoupled, mode \(k\) is governed by
\begin{equation}
  s^2+\delta_k s+\nu_k=0.
  \label{eq:fixed_scalar_generalized}
\end{equation}
The damping ratio is
\begin{equation}
  \zeta_k=\frac{\delta_k}{2\sqrt{\nu_k}}.
  \label{eq:fixed_zeta_generalized}
\end{equation}
This formula is a screen outside the decoupled case. It is still extremely
useful because it exposes the competing effects:
\begin{equation}
  \frac{d\zeta_k}{\zeta_k}
  =
  \frac{d\delta_k}{\delta_k}
  -
  \frac{1}{2}\frac{d\nu_k}{\nu_k}.
  \label{eq:zeta_log_differential}
\end{equation}
An action improves damping ratio only if it raises modal damping faster than
half the relative increase in modal stiffness.

\section{Uniform-Decay Case \texorpdfstring{\(D=\gamma M\)}{D=gamma M}}

The most restrictive decoupled case is
\begin{equation}
  D=\gamma M,\qquad \gamma>0.
  \label{eq:proportional_damping}
\end{equation}
Then \(\delta_k=\gamma\), \(\Gamma=\gamma I\), and every nonzero graph mode
satisfies
\begin{equation}
  s^2+\gamma s+\nu_k=0.
  \label{eq:fixed_scalar_mode}
\end{equation}
The roots are
\begin{equation}
  s_{k,\pm}=\frac{-\gamma\pm\sqrt{\gamma^2-4\nu_k}}{2}.
  \label{eq:fixed_proportional_roots}
\end{equation}
When \(\nu_k>\gamma^2/4\), the roots are underdamped and
\begin{equation}
  \Real(s_{k,\pm})=-\gamma/2.
  \label{eq:fixed_decay_floor}
\end{equation}

\begin{proposition}[Proportional-damping floor]
Under \(D=\gamma M\), every underdamped nonzero mode has the same decay rate
\(-\gamma/2\). Network topology changes \(\nu_k\) and therefore modal
frequency, but it does not change the real decay rate of an underdamped mode
inside the proportional class.
\end{proposition}

\begin{proof}
Equation~\eqref{eq:fixed_scalar_mode} gives the two roots directly. If
\(\nu_k>\gamma^2/4\), the square root is imaginary and the real part is
\(-\gamma/2\). The expression contains no dependence on \(\nu_k\) in its real
part.
\end{proof}

The proposition is sometimes counterintuitive because it does not say topology
is irrelevant. Topology can move a mode between overdamped and underdamped
regions, and it changes frequency. It states only that, after a mode is
underdamped, its decay rate is fixed by \(\gamma\) in this uniform
damping-to-inertia model.

\section{Commuting and Generalized Proportional Damping}

Uniform damping is not the whole decouplable class. The exact condition is
\begin{equation}
  [\Lt,\Dt]=0.
  \label{eq:fixed_commuting_decoupling}
\end{equation}
For real symmetric matrices, \eqref{eq:fixed_commuting_decoupling} is
equivalent to simultaneous orthogonal diagonalization. Therefore there exists
an orthogonal \(Q\) such that
\begin{equation}
  Q^T\Lt Q=\diag(\nu_k),
  \qquad
  Q^T\Dt Q=\diag(\delta_k).
  \label{eq:fixed_simultaneous_diagonalization}
\end{equation}
Each mode then satisfies
\begin{equation}
  s_k^2+\delta_k s_k+\nu_k=0,
  \qquad
  \zeta_k=\frac{\delta_k}{2\sqrt{\nu_k}}.
  \label{eq:fixed_commuting_scalar_modes}
\end{equation}

Rayleigh damping is the most important generalized proportional example:
\begin{equation}
  D=\alpha M+\beta L.
  \label{eq:rayleigh_damping}
\end{equation}
In normalized coordinates,
\begin{equation}
  \Dt=\alpha I+\beta\Lt,
  \label{eq:rayleigh_normalized}
\end{equation}
so \([\Lt,\Dt]=0\). Mode \(k\) satisfies
\begin{equation}
  s_k^2+(\alpha+\beta\nu_k)s_k+\nu_k=0,
  \label{eq:rayleigh_mode_polynomial}
\end{equation}
and its underdamped real part and damping ratio are
\begin{align}
  \Real(s_k)&=-\frac{\alpha+\beta\nu_k}{2},
  \label{eq:rayleigh_real_part}\\
  \zeta_k&=
  \frac{\alpha}{2\sqrt{\nu_k}}
  +
  \frac{\beta\sqrt{\nu_k}}{2}.
  \label{eq:rayleigh_zeta}
\end{align}
Unlike the \(D=\gamma M\) floor, Rayleigh damping lets network stiffness change
both frequency and decay. Depending on the modal range and on the signs and
magnitudes of \(\alpha\) and \(\beta\), reinforcement can improve one set of
modes while degrading another damping-ratio target. This is why the thesis
uses the commuting condition as the decoupling test and reserves
\(D=\gamma M\) for the uniform-decay case.

\section{Exact QEP Sensitivities}

Let a scalar action \(\rho\) perturb all three operators:
\begin{equation}
  P(s,\rho)=s^2M(\rho)+sD(\rho)+L(\rho).
\end{equation}
For a simple eigenvalue \(s_\star\), right vector \(v\), and left vector \(w\),
the QEP sensitivity is
\begin{equation}
  \frac{ds_\star}{d\rho}
  =
  -\frac{
  w^H\left(s_\star^2M_\rho+s_\star D_\rho+L_\rho\right)v
  }{
  w^H\left(2s_\star M+D\right)v
  }.
  \label{eq:fixed_full_qep_sens}
\end{equation}
This formula, standard in QEP perturbation theory \cite{tisseur2001qep}, is
the most compact answer to how \(d/dL\), \(d/dD\), and \(d/dM\) mix.

The three numerator terms have different phase factors:
\begin{equation}
  L_\rho,\qquad s_\star D_\rho,\qquad s_\star^2M_\rho.
  \label{eq:phase_weighted_terms}
\end{equation}
For an oscillatory pole, those terms are not collinear in the complex plane.
Thus, line reinforcement, damping addition, and inertia addition do not point
in the same eigenvalue direction. A scalar ranking that ignores these phase
weights is not a damping certificate.

For \(s=\sigma+j\omega\), the damping-ratio derivative is
\begin{equation}
  \frac{d\zeta}{d\rho}
  =
  -\frac{\omega^2}{(\sigma^2+\omega^2)^{3/2}}
  \Real\left(\frac{ds}{d\rho}\right)
  +
  \frac{\sigma\omega}{(\sigma^2+\omega^2)^{3/2}}
  \Imag\left(\frac{ds}{d\rho}\right).
  \label{eq:fixed_zeta_sensitivity}
\end{equation}
The real-part movement alone is insufficient when the frequency movement is
large.

\section{Derivative with Respect to the Graph \texorpdfstring{\(L\)}{L}}

Let a candidate line \(e=(i,j)\) have incidence vector \(a_e=e_i-e_j\) and
susceptance-like weight \(b_e\). In the Laplacian approximation,
\begin{equation}
  L=\sum_{e} b_e a_ea_e^T,\qquad
  \frac{\partial L}{\partial b_e}=a_ea_e^T.
  \label{eq:laplacian_edge_sum}
\end{equation}
For the generalized graph mode \eqref{eq:generalized_modes},
\begin{equation}
  \frac{\partial \nu_k}{\partial b_e}
  =
  (\phi_k(i)-\phi_k(j))^2.
  \label{eq:edge_spectral_derivative}
\end{equation}
Equation~\eqref{eq:edge_spectral_derivative} is the core weak-link screen:
edges across large modal angle differences have the strongest effect on modal
frequency.

For the scalar decoupled damping-ratio screen with fixed \(D\),
\begin{equation}
  \frac{\partial \zeta_k}{\partial b_e}
  =
  -\frac{\zeta_k}{2\nu_k}
  (\phi_k(i)-\phi_k(j))^2
  +
  \frac{\zeta_k}{\delta_k}
  \frac{\partial \delta_k}{\partial b_e}.
  \label{eq:zeta_edge_derivative_general}
\end{equation}
If modal damping is frozen under the line action, the second term is zero and
line reinforcement reduces \(\zeta_k\) for that mode. It may still improve
angle coherence, synchronizing torque, or protected resolvent gain, but it does
not improve scalar damping ratio by itself.

\begin{example}[A reinforced cut can hide a damping loss]
Suppose a critical mode has \(\delta_k\) fixed and a line across its modal cut
has \((\phi_k(i)-\phi_k(j))^2\) large. Reinforcing that line raises
\(\nu_k\), which is good for graph stiffness. Yet
\eqref{eq:zeta_edge_derivative_general} gives
\(\partial\zeta_k/\partial b_e<0\) if \(\partial\delta_k/\partial b_e=0\).
The graph has become stiffer, but the damping ratio has decreased. This is a
hidden margin failure.
\end{example}

\section{Derivative with Respect to Damping \texorpdfstring{\(D\)}{D}}

For a local damping addition at bus \(i\),
\begin{equation}
  D_\rho=e_ie_i^T.
\end{equation}
The frozen-mode damping derivative is
\begin{equation}
  \frac{\partial\delta_k}{\partial D_i}=\phi_k(i)^2.
  \label{eq:damping_node_derivative}
\end{equation}
Therefore
\begin{equation}
  \frac{\partial\zeta_k}{\partial D_i}
  =
  \frac{\phi_k(i)^2}{2\sqrt{\nu_k}}
  \label{eq:zeta_damping_derivative}
\end{equation}
in the decoupled scalar screen. The best damping location for a single mode is
an antinode of that mode, not necessarily the largest generator, the largest
load, or the highest-degree bus.

For multiple modes, a risk-weighted damping placement score is
\begin{equation}
  \mathcal D_i
  =
  \sum_{k\in\mathcal K}
  \omega_k
  \frac{\phi_k(i)^2}{2\sqrt{\nu_k}},
  \label{eq:damping_placement_score}
\end{equation}
where \(\mathcal K\) is the set of monitored modes and \(\omega_k\ge0\) weights
criticality. This is a screen. The exact acceptance test remains
\eqref{eq:fixed_full_qep_sens} or full eigenvalue recomputation.

\begin{example}[A low-degree bus can be the best damping site]
A bus with modest graph degree can have large \(\phi_k(i)^2\) for the critical
mode. Equation~\eqref{eq:zeta_damping_derivative} ranks it highly for damping
placement even if graph centrality ranks it poorly. This is a weak-node
mechanism: the bus is weak because the critical mode lives there.
\end{example}

\section{Derivative with Respect to Inertia \texorpdfstring{\(M\)}{M}}

For a local inertia addition at bus \(i\),
\begin{equation}
  M_\rho=e_ie_i^T.
\end{equation}
The generalized eigenvalue derivative is
\begin{equation}
  \frac{\partial\nu_k}{\partial M_i}
  =
  -\nu_k\phi_k(i)^2.
  \label{eq:inertia_frequency_derivative}
\end{equation}
Local inertia lowers the modal frequency most strongly where the mode has large
amplitude. That fact supports inertia placement for frequency nadir or RoCoF
objectives, and it is consistent with studies showing that virtual-inertia
placement can be posed as a network performance problem
\cite{poolla2017virtual}, that inertia location and slow modes shape
disturbance propagation \cite{pagnier2019inertia_location}, and that inertia
and primary-control placement separate under heterogeneous parameters
\cite{pagnier2019optimal}.

For oscillatory damping ratio, however, the conclusion is more delicate. The
exact answer is still the QEP derivative \eqref{eq:fixed_full_qep_sens}. As a
fixed-mode planning screen, if physical damping \(D\) is not co-added, the
frozen-shape normalization gives
\begin{equation}
  \frac{\partial\delta_k}{\partial M_i}
  \approx
  -\delta_k\phi_k(i)^2.
  \label{eq:inertia_damping_dilution}
\end{equation}
Combining \eqref{eq:zeta_log_differential},
\eqref{eq:inertia_frequency_derivative}, and
\eqref{eq:inertia_damping_dilution} yields
\begin{equation}
  \frac{\partial\zeta_k}{\partial M_i}
  \approx
  -\frac{1}{2}\zeta_k\phi_k(i)^2.
  \label{eq:inertia_zeta_dilution}
\end{equation}
Thus, in a fixed-\(D\) model, inertia-only placement at a modal antinode can
decrease the scalar damping-ratio screen. It may still improve RoCoF or reduce
frequency peaks, but it is not a free damping improvement. Any candidate that
matters operationally should therefore be rechecked with
\eqref{eq:fixed_full_qep_sens} or with a full eigenvalue recomputation.

If inertia is co-designed with proportional damping,
\begin{equation}
  D_i=\gamma M_i,
  \label{eq:co_designed_proportional}
\end{equation}
then the damping-to-inertia ratio is preserved. The real decay floor remains
\(-\gamma/2\), and the damping ratio increases as \(\nu_k\) decreases. This is
the engineering difference between adding inertia alone and adding a
properly-damped inertial device.

\begin{example}[Inertia placement has two different objectives]
Placing inertia where \(\phi_k(i)^2\) is large is attractive if the objective
is to slow a disturbance. Under fixed \(D\), the same placement can reduce
\(\zeta_k\) through damping dilution. A correct planning statement must say
whether the target is RoCoF, frequency nadir, decay rate, damping ratio, or
input-output amplification.
\end{example}

\section{Graph Robustness Without Adding Inertia}

If inertia is unavailable or undesirable, the graph can be strengthened through
line weights, switching, topology control, or series compensation. The oldest
spectral screen is algebraic connectivity: maximize \(\nu_2\), the first
nonzero Laplacian eigenvalue introduced by Fiedler's graph connectivity theory
\cite{fiedler1973algebraic}. The local derivative is
\begin{equation}
  \frac{\partial\nu_2}{\partial b_e}
  =
  (\phi_2(i)-\phi_2(j))^2.
  \label{eq:fiedler_edge_score}
\end{equation}
This favors lines across the Fiedler cut.

A different robustness screen is total effective resistance, also called the
Kirchhoff index:
\begin{equation}
  K_f=n\,\mathrm{trace}(L^+)
  =
  n\sum_{k=2}^n\frac{1}{\lambda_k}.
  \label{eq:kirchhoff_index}
\end{equation}
Minimizing effective resistance is a convex edge-weight allocation problem in
standard graph settings \cite{ghosh2008effective}. Its derivative is
\begin{equation}
  \frac{\partial\,\mathrm{trace}(L^+)}{\partial b_e}
  =
  -a_e^T(L^+)^2a_e
  =
  -\sum_{k=2}^n
  \frac{(q_k(i)-q_k(j))^2}{\lambda_k^2}.
  \label{eq:effective_resistance_derivative}
\end{equation}
This favors lines that reduce many slow modes, not only the Fiedler mode.
Effective resistance also underlies spectral sparsification, where edges are
sampled or weighted according to their resistance contribution while preserving
quadratic forms \cite{spielman2011sparsification}.

The hidden damping question remains:
\begin{equation}
  \hbox{Does the graph action also improve }
  \zeta_k,\ \Real(s_k),\hbox{ or }g_{\rm prot}?
  \label{eq:graph_hidden_question}
\end{equation}
If not, the action is graph-robust but damping-ambiguous. It must be paired
with damping, controller retuning, or a certificate that the chosen margin is
not damping ratio.

\section{Weak Nodes, Weak Links, and Hidden Margins}

The fixed-\(D\) theory gives precise first-order definitions.

\begin{definition}[Weak node for damping]
For a monitored mode set \(\mathcal K\), bus \(i\) is weak for damping if its
modal damping placement score \(\mathcal D_i\) in
\eqref{eq:damping_placement_score} is large while the current local damping
contribution \(D_i\phi_k(i)^2\) is small for critical modes.
\end{definition}

\begin{definition}[Weak link for graph stiffness]
Line \(e\) is weak for a graph metric if the derivative of that metric with
respect to \(b_e\) is large. Examples are
\eqref{eq:fiedler_edge_score} for algebraic connectivity and
\eqref{eq:effective_resistance_derivative} for total effective resistance.
\end{definition}

\begin{definition}[Hidden damping margin]
A graph action has hidden damping risk for mode \(k\) if
\begin{equation}
  \frac{\partial \nu_k}{\partial\rho}>0
  \quad\hbox{but}\quad
  \frac{\partial \zeta_k}{\partial\rho}<0
  \label{eq:hidden_margin_definition}
\end{equation}
under the selected damping model. It improves a stiffness screen while
degrading a damping-ratio screen.
\end{definition}

This definition prevents a common overclaim: a stronger graph is not the same
as a better damping margin.

\section{A Fixed-\texorpdfstring{\(D\)}{D} Spectral Planning Screen}

The following screen extracts the maximum information available before moving
to the full QEP or Schur NEP.

\begin{enumerate}
\item Compute generalized modes \(L\phi_k=\nu_kM\phi_k\).
\item Compute modal damping \(\delta_k=\phi_k^TD\phi_k\) and
\(\zeta_k=\delta_k/(2\sqrt{\nu_k})\).
\item Compute node damping scores \(\mathcal D_i\).
\item Compute inertia-dilution scores
\begin{equation}
  \mathcal M_i
  =
  \sum_{k\in\mathcal K}
  \omega_k\zeta_k\phi_k(i)^2.
  \label{eq:inertia_dilution_score}
\end{equation}
\item Compute line graph scores using \(\partial\nu_2/\partial b_e\) and
\(\partial\mathrm{trace}(L^+)/\partial b_e\).
\item Flag hidden-margin conflicts where graph scores are positive but
\(\partial\zeta_k/\partial b_e<0\).
\item Verify accepted candidates with the exact QEP derivative
\eqref{eq:fixed_full_qep_sens}.
\end{enumerate}

This screen is intentionally not the final thesis method. It is the strongest
fixed-\(D\) layer before controller self-energy is introduced.
